\chapter{\label{marcoteorico}Marco Teórico}

\section{Contexto y Robótica de Servicio}
La robótica de servicio se orienta a la asistencia en tareas cotidianas en entornos no industriales, priorizando la interacción flexible con seres humanos. En el ámbito de los UAVs, el seguimiento autónomo es una extensión de esta disciplina, con aplicaciones en educación, seguridad y asistencia personalizada.

\section{Tecnologías de Detección y Seguimiento}
El seguimiento visual en tiempo real requiere la combinación de visión por computador y aprendizaje profundo (Deep Learning). 

\begin{itemize}
    \item \textbf{Redes Neuronales Convolucionales (CNN):} Algoritmos avanzados diseñados para el procesamiento de imágenes y reconocimiento visual.
    \item \textbf{Arquitecturas Ligeras:} Modelos como \textbf{MobileNetV2} son fundamentales para sistemas embebidos debido a su estructura basada en convoluciones separables.
\end{itemize}

\section{Software y Middleware: ROS (Robot Operating System)}
El uso de ROS permite una creación modular de aplicaciones robóticas complejas mediante comunicación por nodos y tópicos.
